% GENERATED FILE. Edit canonical lessons, then run npm run generate:books.
% canonical-source-hash: fnv1a64:e9793a24711b27a9
% canonical-lessons: ZH-C436-dixing, ZH-C436-nanfang, ZH-W436-fang-side, ZH-C436-beifang, ZH-C436-fangyan, ZH-C436-kouyin

\chapter{North and South}
\label{ch:zh-north-and-south}
\hlchaptermodality{436}

\begin{chapteropening}
\textbf{By the end of this chapter:} \emph{I can talk about the lie of the land, the south and the north, a dialect and an accent.}

Complete the last lesson of chapter 436: I can talk about the lie of the land, the south and the north, a dialect and an accent.
\end{chapteropening}

\section[dìxíng]{dìxíng — terrain; lie of the land}

\label{lesson:ZH-C436-dixing}



\begin{quote}
Before the new one: say the Mandarin for spectacular, then the Mandarin for to come down a mountain.
\end{quote}

\subsection*{You'll want to know: dìxíng}
\emph{dìxíng} — \textquotedblleft{}terrain; lie of the land\textquotedblright{}.

\begin{quote}
\emph{zhège dìqū de dìxíng hěn tèbié} — The land in this region has an unusual shape.
\end{quote}

Its characters come later; for now it is a word for the ear.

\begin{cousinweb}
\emph{dìxíng} is built from \zh{地} 'the earth' and \emph{xíng} 'a shape', both inherited from Middle Chinese, \emph{dijH} and \emph{heng} (Guangyun, in Baxter's transcription), which glosses the second as the look of a thing.
\end{cousinweb}

\begin{grammarlens}[title={What you've built}]
One more word for getting around.
\end{grammarlens}

\subsection*{Your turn}
\begin{itemize}
\raggedright
  \item \emph{Say it:} \emph{dìxíng}
  \item \emph{Say it:} \emph{dìxíng}, once more
  \item \emph{Say it:} \emph{zhège dìqū de dìxíng hěn tèbié} — The land in this region has an unusual shape.
  \item \emph{Recall:} say the Mandarin for spectacular, then the Mandarin for to come down a mountain, then say \emph{dìxíng} again
\end{itemize}

\begin{tcolorbox}[breakable,colback=teal!4,colframe=teal!35!black,title={Before you move on}]
What does dìxíng mean? (\textquotedblleft{}Terrain; lie of the land\textquotedblright{}.) Say it once more.
\end{tcolorbox}
\section[\texorpdfstring{\zh{南方} (nánfāng)}{南方 (nánfāng)}]{\zh{南方} (nánfāng) — the south of the country; southern China}

\label{lesson:ZH-C436-nanfang}



\begin{quote}
Before the new one: say the Mandarin for to come down a mountain, then the Mandarin for the lie of the land.
\end{quote}

\subsection*{You'll want to know: \zh{南方}}
\textbf{\zh{南方}} — \emph{nánfāng} — \textquotedblleft{}the south of the country; southern China\textquotedblright{}.

\begin{quote}
\emph{dōngtiān wǒmen zuò huǒchē qù nánfāng} — In winter we take the train down south.
\end{quote}

The new character, in pieces: \zh{方} — four strokes: a dot, a long bar, a stroke that turns down and hooks, and a long stroke falling to the left. The next lesson puts it in your hand.

\begin{cousinweb}
\textbf{\zh{南方}} (\emph{nánfāng}) is built from \zh{南} 'south' and \zh{方} 'a quarter', both inherited from Middle Chinese, \emph{nom} and \emph{pjang} (Guangyun, in Baxter's transcription), which glosses the second as the four quarters of the world.
\end{cousinweb}

\begin{grammarlens}[title={What you've built}]
One more word for getting around.
\end{grammarlens}

\subsection*{Your turn}
\begin{itemize}
\raggedright
  \item \emph{Say it:} \emph{nánfāng}
  \item \emph{Say it:} \emph{nánfāng}, once more
  \item \emph{Say it:} \emph{dōngtiān wǒmen zuò huǒchē qù nánfāng} — In winter we take the train down south.
  \item \emph{Recall:} say the Mandarin for to come down a mountain, then the Mandarin for the lie of the land, then say \emph{nánfāng} again
\end{itemize}

\begin{tcolorbox}[breakable,colback=teal!4,colframe=teal!35!black,title={Before you move on}]
What does \zh{南方} mean? (\textquotedblleft{}The south of the country; southern China\textquotedblright{}.) Say it once more.
\end{tcolorbox}
\section[\texorpdfstring{\zh{方} (fāng)}{方 (fāng)}]{\zh{方} — a dot over a bar, a hook and a long stroke}

\label{lesson:ZH-W436-fang-side}



\begin{quote}
The word first, then the character inside it.

\begin{itemize}
\raggedright
  \item \emph{Recall:} say \emph{nánfāng} — \textbf{R1}, one lesson back, before it decays
\end{itemize}

\begin{quote}
\zh{南} \zh{万}
\end{quote}

\emph{Read it:} each one aloud, with nothing said first — \emph{nán}, \emph{wàn}
\end{quote}

\subsection*{Script — a dot over a bar, a hook and a long stroke}
\hlblockfigure{\detokenize{figures/ZH-W436-fang-side-filmstrip.pdf}}{How \zh{方} is written, stroke by stroke}

\begin{quote}
\zh{方}
\end{quote}

\textbf{One piece, four strokes.} Count the strokes this time: there is only one piece.

\textbf{A dot on top}, in the middle.

\textbf{A long bar under it.}

\textbf{From under the bar, a stroke that runs a little to the right, turns down and hooks to the left at the foot.} Then \textbf{a long stroke falling to the left}, also from under the bar. These last two are made as in \zh{万}.

1 + 1 + 2 = four.

\subsection*{Your turn}
\begin{itemize}
\raggedright
  \item \emph{Write it:} \zh{方}, all four, with the model in front of you
  \item \emph{Cover:} write it again from memory, and count the pieces, not the strokes
\end{itemize}

\begin{tcolorbox}[breakable,colback=teal!4,colframe=teal!35!black,title={Before you move on}]
How many strokes in \zh{方}? (\textbf{Four}.) What comes first? (\textbf{A dot}.) Which stroke hooks? (\textbf{The third: it runs right, turns down and hooks}.)

Source: \href{https://raw.githubusercontent.com/chanind/hanzi-writer-data/68d10a4b21150cae5e1ebbd223eed289cf32d90c/data/\%E6\%96\%B9.json}{Hanzi Writer Data \zh{方}, pinned stroke paths and medians}.
\end{tcolorbox}
\section[\texorpdfstring{\zh{北方} (běifāng)}{北方 (běifāng)}]{\zh{北方} (běifāng) — the north of the country; northern China}

\label{lesson:ZH-C436-beifang}



\begin{quote}
Before the new one: say the Mandarin for the lie of the land, then the Mandarin for the south of the country.
\end{quote}

\subsection*{You'll want to know: \zh{北方}}
\textbf{\zh{北方}} — \emph{běifāng} — \textquotedblleft{}the north of the country; northern China\textquotedblright{}.

\begin{quote}
\emph{wǒ zhù zài běifāng dōngtiān hěn lěng} — I live in the north, where the winters are cold.
\end{quote}

\begin{cousinweb}
\textbf{\zh{北方}} (\emph{běifāng}) is built from \zh{北} 'north' and \zh{方} 'a quarter', both inherited from Middle Chinese, \emph{pok} and \emph{pjang} (Guangyun, in Baxter's transcription).
\end{cousinweb}

\begin{grammarlens}[title={What you've built}]
One more word for getting around.
\end{grammarlens}

\subsection*{Your turn}
\begin{itemize}
\raggedright
  \item \emph{Say it:} \emph{běifāng}
  \item \emph{Say it:} \emph{běifāng}, once more
  \item \emph{Say it:} \emph{wǒ zhù zài běifāng dōngtiān hěn lěng} — I live in the north, where the winters are cold.
  \item \emph{Recall:} say the Mandarin for the lie of the land, then the Mandarin for the south of the country, then say \emph{běifāng} again
\end{itemize}

\begin{tcolorbox}[breakable,colback=teal!4,colframe=teal!35!black,title={Before you move on}]
What does \zh{北方} mean? (\textquotedblleft{}The north of the country; northern China\textquotedblright{}.) Say it once more.
\end{tcolorbox}
\section[fāngyán]{fāngyán — dialect}

\label{lesson:ZH-C436-fangyan}



\begin{quote}
Before the new one: say the Mandarin for the south of the country, then the Mandarin for the north of the country.
\end{quote}

\subsection*{You'll want to know: fāngyán}
\emph{fāngyán} — \textquotedblleft{}dialect\textquotedblright{}.

\begin{quote}
\emph{zhèr de rén shuō fāngyán wǒ tīng bu dǒng} — People here speak a dialect I can't understand.
\end{quote}

Its characters come later; for now it is a word for the ear.

\begin{cousinweb}
\emph{fāngyán} is built from \zh{方} 'a region' and \emph{yán} 'speech', both inherited from Middle Chinese, \emph{pjang} and \emph{ngjon} (Guangyun, in Baxter's transcription), which glosses the second 'speech'. The Guangyun itself quotes an old book called the Fangyan, a record of the speech of the regions.
\end{cousinweb}

\begin{grammarlens}[title={What you've built}]
One more word for getting around.
\end{grammarlens}

\subsection*{Your turn}
\begin{itemize}
\raggedright
  \item \emph{Say it:} \emph{fāngyán}
  \item \emph{Say it:} \emph{fāngyán}, once more
  \item \emph{Say it:} \emph{zhèr de rén shuō fāngyán wǒ tīng bu dǒng} — People here speak a dialect I can't understand.
  \item \emph{Recall:} say the Mandarin for the south of the country, then the Mandarin for the north of the country, then say \emph{fāngyán} again
\end{itemize}

\begin{tcolorbox}[breakable,colback=teal!4,colframe=teal!35!black,title={Before you move on}]
What does fāngyán mean? (\textquotedblleft{}Dialect\textquotedblright{}.) Say it once more.
\end{tcolorbox}
\section[\texorpdfstring{\zh{口音} (kǒuyīn)}{口音 (kǒuyīn)}]{\zh{口音} (kǒuyīn) — accent}

\label{lesson:ZH-C436-kouyin}



\begin{quote}
Before the new one: say the Mandarin for the north of the country, then the Mandarin for a regional dialect.
\end{quote}

\subsection*{You'll want to know: \zh{口音}}
\textbf{\zh{口音}} — \emph{kǒuyīn} — \textquotedblleft{}accent\textquotedblright{}.

\begin{quote}
\textbf{\zh{他说话有南方口音}} \emph{tā shuōhuà yǒu nánfāng kǒuyīn} — He speaks with a southern accent.
\end{quote}

\begin{cousinweb}
\textbf{\zh{口音}} (\emph{kǒuyīn}) is built from \zh{口} 'the mouth' and \zh{音} 'a sound', both inherited from Middle Chinese, \emph{khuwX} and \emph{'im} (Guangyun, in Baxter's transcription), which glosses the second 'a sound'.
\end{cousinweb}

\begin{grammarlens}[title={What you've built}]
One more word for getting around.
\end{grammarlens}

\subsection*{Your turn}
\begin{itemize}
\raggedright
  \item \emph{Say it:} \emph{kǒuyīn}
  \item \emph{Say it:} \emph{kǒuyīn}, once more
  \item \emph{Say it:} \emph{tā shuōhuà yǒu nánfāng kǒuyīn} — He speaks with a southern accent.
  \item \emph{Recall:} say the Mandarin for the north of the country, then the Mandarin for a regional dialect, then say \emph{kǒuyīn} again
\end{itemize}

\begin{tcolorbox}[breakable,colback=teal!4,colframe=teal!35!black,title={Before you move on}]
What does \zh{口音} mean? (\textquotedblleft{}Accent\textquotedblright{}.) Say it once more.
\end{tcolorbox}
